\documentclass[11pt]{article}

\usepackage[margin=1in]{geometry}
\usepackage{hyperref}
\usepackage{graphicx}

\title{DIP-001: Decision Integrity Protocol (DIP)\\
A Cryptographic Framework for Deterministic Governance}

\author{Pavan Dev Singh Charak}

\date{2026}

\begin{document}

\maketitle

\begin{abstract}

The Decision Integrity Protocol (DIP) introduces a cryptographic
framework for deterministic governance and verifiable decision systems.
The protocol ensures that organizational and automated decisions
produce verifiable evidence that can be independently audited.

\end{abstract}

\section{Introduction}

Modern organizations rely heavily on automated systems and complex
decision processes. However, many systems lack mechanisms to produce
verifiable evidence describing how decisions were made.

The Decision Integrity Protocol addresses this challenge by introducing
a deterministic framework for documenting and verifying decision events.

\section{Problem Statement}

Decision systems often suffer from:

\begin{itemize}
\item lack of verifiable decision history
\item fragmented audit trails
\item inability to reproduce decisions
\end{itemize}

These limitations create governance and accountability challenges.

\section{Protocol Architecture}

The Decision Integrity Protocol consists of four primary components:

\begin{itemize}
\item Documentation Engine
\item Decision Payload
\item Decision Ledger
\item Verification Layer
\end{itemize}

These components ensure that every decision produces cryptographically
verifiable evidence.

\section{Decision Payload}

A decision payload represents the canonical record of a decision event.

\begin{verbatim}
{
 "decision_id": "abc123",
 "timestamp": "2026-03-06T10:30:00Z",
 "actor": "system",
 "action": "approve_transaction",
 "result": "approved"
}
\end{verbatim}

\section{Decision Ledger}

The decision ledger is an append-only registry that stores signed
decision payloads.

This ensures that historical decision records cannot be altered.

\section{Verification}

Independent verifiers can validate decision integrity by checking:

\begin{itemize}
\item payload canonicalization
\item signature validity
\item ledger inclusion
\end{itemize}

\section{Conclusion}

Decision Integrity Protocol enables deterministic governance systems
that provide cryptographically verifiable decision records.

\end{document}